We ask how much of a knowledge-graph completion model's accuracy survives a move from held-out triples of its training graph to a freshly generated graph with entirely new entities, in a setting where the relations are derived by rules and the answer structure is known by construction. We generate family-tree graphs of \rhval{const/people} people with four relations (parent, sibling, grandparent, uncle/aunt), and compare a tuned, reimplemented TransE, TransE with entity fold-in, a 2-layer query-conditioned relational message-passing model (Path-MP, a reduced NBFNet-style network) and a hand-written one-step rule oracle, over \rhval{const/n_seeds} seeds that each use a different graph pair. Path-MP reaches filtered MRR \rhval{main/path-mp-2-hop/transductive/mrr/mean:3} on held-out triples and \rhval{main/path-mp-2-hop/inductive/mrr/mean:3} on the new graph, with no measurable drop; TransE reaches \rhval{main/transe-reimplemented/transductive/mrr/mean:3} and falls to \rhval{main/transe-reimplemented/inductive/mrr/mean:3}. Fitting new entity vectors with frozen relation vectors restores TransE to \rhval{main/transe+fold-in-reimplemented/inductive/mrr/mean:3}, still far below Path-MP. Ablations show that one message-passing layer is at chance, a third layer adds nothing measurable, training without removing the query edge costs about half the MRR, and when only a fraction of the new graph is observed Path-MP degrades quickly (MRR \rhval{analysis/analysis-final/inductive/dens_pm_f25_mean/mean:3} at a quarter of the triples) and is no longer distinguishable from fold-in TransE. The study is small, synthetic and has a high ceiling for the path model; it measures a gap that is known in principle, and we report where it is and is not there.
